\documentclass{report}
\usepackage{amsmath,amssymb}
\setlength{\parindent}{0mm}
\setlength{\parskip}{1em}
\begin{document}
\begin{center}
\rule{6in}{1pt} \
{\large Tom Jonsthovel \\
{\bf A comparison of a rigid body mode deflation method and AMG for simulation of 3D mechanics}}

Schlumberger Technology Center \\ Lambourn Court \\ Wyndyke Furlong \\ OX14 1UJ \\ Abingdon
\\
{\tt tjonsthovel@slb.com}\\
Scott MacLachlan\\
Martin van Gijzen\end{center}

In many real life engineering problems the analysis of mechanical systems
by simulation plays an important role. Moreover, studying the behavior of
these mechanical systems is important as they are often coupled to other
physical phenomena, an example is geomechanics and fluid flow in
reservoir engineering. Simulation of mechanical problems involves
modeling of materials with non-linear material properties like hyper
elasticity, plasticity and viscosity. These non-linear equations are
often solved by the Newton-Raphson method which involves solving one
linear system per Newton-Raphson iteration. Solving the linear systems
dominates the total runtime of the simulation and hence has to be fast.
For realistic three dimensional mechanical problems that are defined on
meshes with a large number of elements and that involve different
materials, the resulting stiffness matrix is often sparse and ill
conditioned. We need powerful preconditioners to solve these linear
systems with iterative solution methods. In this talk we compare the
performance of two preconditioners for the conjugate gradient method: the
deflation method based on the elimination of the rigid body modes and
Smoothed Aggregation Algebraic Multigrid (SA AMG). We show and discuss
the results for matrices from two different real life applications,
structural mechanics and geomechanics.


\end{document}
